\chapter{Proposed Work}
\label{sec:proposed-work}

\todo[inline]{Two opening paragraphs: the problem in one sentence; what
the chapter defines.}

\section{Proposed Approach}
\label{sec:proposed-work-approach}
\todo[inline]{The idea in prose: citizen-science observations, mask,
labels, embeddings, classifier; expert data only to measure the gain.}

\section{Reference-Data Pipeline}
\label{sec:proposed-work-pipeline}

\subsection{Citizen-science observations}
\label{sec:proposed-work-pipeline-cs}
\todo[inline]{}

\subsection{Expert polygons}
\label{sec:proposed-work-pipeline-expert}
\todo[inline]{}

\subsection{From points to pixel labels}
\label{sec:proposed-work-pipeline-labels}
\todo[inline]{}

\section{Experimental Design}
\label{sec:proposed-work-design}

\subsection{Task and classifier}
\label{sec:proposed-work-design-task}
\todo[inline]{}

\subsection{Spatial validation}
\label{sec:proposed-work-design-validation}
\todo[inline]{}

\subsection{Metric}
\label{sec:proposed-work-design-metric}
\todo[inline]{}

\subsection{Hypotheses}
\label{sec:proposed-work-design-hypotheses}
\todo[inline]{H$_0$/H$_1$ per research question, as in
docs/research\_questions.md.}

\section{Dataset Description and Preparation}
\label{sec:proposed-work-data}

This section describes the study area and each dataset as it was
obtained during the first semester: the iNaturalist observations, the
expert polygons, the TESSERA embeddings, and the state of the
Sentinel-2 and land-cover data. The steps that turn the observations
and polygons into labelled pixels are those of the reference-data
pipeline (Section~\ref{sec:proposed-work-pipeline}).

\subsection{Study area}
\label{sec:proposed-work-data-area}

The study area is the Coimbra district, in central Portugal. It is the
area of the iNaturalist place used to query the observations (place
8422, ``Coimbra, PT''), which resolves to the district and not to the
municipality. Sentinel-2 tile 29TNE, on which most expert polygons
were drawn, covers the district. All data are handled on the
10~m grid of the TESSERA embeddings, in UTM zone 29N (EPSG:32629). The
sample of the landscape used by the classifier and the metric is
drawn within a rectangle of 100 by 66~km over the district.
\todo{Study-area map (district, presences, expert polygons) and a
short land-cover description; neither is made yet.}

\subsection{iNaturalist}
\label{sec:proposed-work-data-inaturalist}

The citizen-science observations were downloaded from the iNaturalist
API on 22 September 2026, for genus \emph{Acacia} and all its
descendants within place 8422, with no filter on date or quality grade.
The download holds 7\,943 observations from 1992 to 2026, each with
its date, coordinates, positional accuracy, quality grade, observer
and licence; the licence is set per observation by its author.
Filtering is left to the pipeline, which keeps 2\,465
\emph{A.~dealbata} presences, one per 10~m pixel, observed between
2021 and 2025 (Section~\ref{sec:proposed-work-pipeline-cs}).

\subsection{Expert reference data}
\label{sec:proposed-work-data-expert}

The expert reference data are polygons of \emph{A.~dealbata} stands,
provided by the supervisor as project data. They were made from field
visits and from visual interpretation of OrtoSat 2023 orthoimagery at
30~cm, and refer to 2023. They come in two files: 7\,057 polygons on
Sentinel-2 tile 29TNE and 437 elsewhere in continental Portugal. Each
polygon has a single attribute, whether \emph{A.~dealbata} is present
(1) or not (0). Of the 2\,470 presence polygons, 1\,849 lie in the
Coimbra district. The 5\,013 polygons marked as not \emph{A.~dealbata}
are not used (Section~\ref{sec:proposed-work-pipeline-expert}).

\subsection{TESSERA embeddings}
\label{sec:proposed-work-data-tessera}

The embeddings are those of TESSERA version 1.1 for 2023, read with the
\texttt{geotessera} library (version 0.11.0) at the centre of each
pixel of interest, without downloading whole tiles. Each pixel has 128
values. Two sets were read. The first holds the reference pixels,
about 29\,000 points: the pixels of the citizen-science observations
and of the expert polygons, and points of an earlier design that are
not used; 2\,464 of the 2\,465
presences have a valid embedding. The second is the sample of the
landscape: 2\,000 points drawn uniformly at random within the
rectangle of the study area (seed 42), excluding a 50~m buffer around
every presence of either source, all with a valid embedding. How this
sample enters training and evaluation is part of the experimental
design (Section~\ref{sec:proposed-work-design}).

\subsection{Sentinel-2 and COS}
\label{sec:proposed-work-data-sentinel2}

The Sentinel-2 features are needed only for RQ4 and are not yet
extracted. The access route was tested on a strip of 1.75 by 8.89~km:
the 61 Level-2A scenes of 2023 with less than 20\% cloud cover were
read from the AWS Earth Search catalogue and resampled onto the
TESSERA grid, so that both representations describe the same pixels.
Extraction over the study area is planned for the second semester
(Section~\ref{sec:proposed-work-plan-second}).

The land-cover map COS 2023 (Carta de Uso e Ocupação do Solo) was
obtained from the Direção-Geral do Território. It is not used to
label pixels: its minimum mapping unit of one hectare does not resolve
scattered presences (Section~\ref{sec:background-iap-mapping}). It
records the land-cover class of each expert polygon, for describing the
reference data.

\section{Preliminary Experiment}
\label{sec:proposed-work-preliminary}
\todo[inline]{RQ1 effort-reduction curve: table, figure, buffer
sensitivity, short reading marked as preliminary. Single run:
results/effort\_reduction\_curve/.}

\section{Work Plan}
\label{sec:proposed-work-plan}

\subsection{First Semester}
\label{sec:proposed-work-plan-first}
\todo[inline]{What was done; Gantt with estimated vs.\ actual.}

\subsection{Second Semester}
\label{sec:proposed-work-plan-second}
\todo[inline]{Per research question; Gantt.}

\section{Risk Assessment}
\label{sec:proposed-work-risks}
\todo[inline]{Table: ID, risk, probability, impact, mitigation.}
